
\documentclass[11pt]{article}
\usepackage{amsmath,amssymb,amsthm,mathtools}
\usepackage{booktabs}
\usepackage[margin=1in]{geometry}

\title{Step 103: Semilocal Metric Off-Diagonal Block Audit}
\author{Six Birds / Formed-Layer Membrane RH Program}
\date{}

\newtheorem{theorem}{Theorem}
\newtheorem{lemma}{Lemma}
\newtheorem{proposition}{Proposition}
\newtheorem{definition}{Definition}
\newtheorem{warning}{Warning}

\begin{document}
\maketitle

\section{Purpose}

Step 101 reduced compactness of the semilocal cross-term obstruction
\[
  \Delta_S=q_{\mathrm{Weil},S}-q_{\mathrm{pos},S}
\]
to the compactness of a transported semilocal projection residual.  Step 102
showed that, after pulling the semilocal Hilbert metric back to the
archimedean model,
\[
  \mathcal R_S=U_SP_SU_S^{-1}-P_\infty
\]
is controlled by the off-diagonal metric block
\[
  A_{MN}=P_\infty A_S(I-P_\infty).
\]
This note audits that block for the concrete finite-place semilocal metric.

\section{Semilocal metric}

For a finite set of places \(S\), with \(\infty\in S\), write
\[
  S_f=S\setminus\{\infty\}.
\]
On the Mellin/Fourier response coordinate \(\xi\), the finite local factors
give
\[
  A_S(\xi)=\prod_{p\in S_f}L_p\left(\frac12-i\xi\right),
  \qquad
  L_p(z)=(1-p^{-z})^{-1}.
\]
Thus the pulled semilocal metric is multiplication by
\[
  a_S(\xi):=|A_S(\xi)|^2
  =
  \prod_{p\in S_f}
  \left|1-p^{-1/2+i\xi}\right|^{-2}.
\]
So, in the common response model,
\[
  A_S=M_{a_S}.
\]

\section{Bohr expansion of the finite-place metric}

Let \(r_p=p^{-1/2}\) and \(\theta_p=\xi\log p\).  The elementary Poisson
expansion gives
\[
  |1-r_pe^{i\theta_p}|^{-2}
  =
  \frac{1}{1-r_p^2}
  \sum_{k\in\mathbb Z} r_p^{|k|}e^{ik\theta_p}.
\]
Therefore, for finite \(S_f\),
\[
  a_S(\xi)
  =
  \sum_{\mathbf k\in\mathbb Z^{S_f}}
  c_{\mathbf k}
  e^{i\xi\ell_{\mathbf k}},
\]
where
\[
  c_{\mathbf k}
  =
  \prod_{p\in S_f}
  \frac{p^{-|\mathbf k_p|/2}}{1-p^{-1}},
  \qquad
  \ell_{\mathbf k}
  =
  \sum_{p\in S_f}\mathbf k_p\log p.
\]
The series is absolutely convergent.

\begin{lemma}[Off-diagonal translation decomposition]
Let \(P=P_\infty\) be the imported archimedean Sonin/prolate projection in
the common response model. Then
\[
  P A_S(I-P)
  =
  \sum_{\mathbf k\ne0}
  c_{\mathbf k}\,
  P\tau_{\ell_{\mathbf k}}(I-P),
\]
where \(\tau_a\) denotes the log-side translation corresponding to
multiplication by \(e^{ia\xi}\) on the Mellin/Fourier side.  The constant
mode drops out because \(P(I-P)=0\).
\end{lemma}

\begin{proof}
The expansion of \(a_S\) is absolutely convergent, so multiplication by
\(a_S\) is the norm-convergent sum of the multiplication operators
\(c_{\mathbf k}M_{e^{i\xi\ell_{\mathbf k}}}\).  Under the inverse
Fourier/Mellin transform \(M_{e^{ia\xi}}\) is the translation \(\tau_a\).
Multiplying on the left by \(P\) and on the right by \(I-P\) gives the
claim.
\end{proof}

\section{Compactness gate}

The previous lemma shows that compactness is not a scalar property of the
local Euler factor.  It is a property of the shifted off-diagonal blocks
\[
  P\tau_{\ell}(I-P).
\]

\begin{proposition}[Compactness reduction]
Assume the block inverse appearing in the Step 102 projection formula is
bounded.  If, for every nonzero Bohr shift \(\ell_{\mathbf k}\), the block
\[
  P_\infty\tau_{\ell_{\mathbf k}}(I-P_\infty)
\]
is compact and the corresponding compact tails are summable against
\(c_{\mathbf k}\), then \(A_{MN}\) is compact.  The analogous statement holds
for Hilbert--Schmidt, trace-class, or Schatten classes.
\end{proposition}

\begin{warning}[Raw local-factor multiplication does not create compactness]
On a non-atomic \(L^2\)-space, multiplication by a nonzero bounded function is
not compact.  The local Euler factor can transport compactness through a
bounded invertible equivalence, but it does not by itself create compactness.
Compactness must come from the Sonin/prolate projection geometry, quotienting,
tail control, or source absorption.
\end{warning}

\section{Hard-cutoff noncompactness model}

The following model explains why the compactness gate is nontrivial.  Let
\(P_H\) be multiplication by the indicator of a half-line or interval
complement in \(L^2(\mathbb R)\), and let \(a\ne0\).  If the crossing region
has positive measure, then
\[
  P_H\tau_a(I-P_H)
\]
contains a partial translation between two positive-measure \(L^2\)-regions.
It is therefore a partial isometry on an infinite-dimensional subspace and is
not compact.

Thus finite-place shifts should be treated as noncompact-risk modes unless
the imported Sonin/prolate projection \(P_\infty\) supplies additional
compact smoothing after the legal quotient.

\section{Verdict}

The semilocal metric block has the concrete form
\[
  A_{MN}
  =
  \sum_{\mathbf k\ne0}c_{\mathbf k}
  P_\infty\tau_{\ell_{\mathbf k}}(I-P_\infty).
\]
The archimedean Connes--Consani residual remains compact-after-quotient, but
the finite-place semilocal metric block is not compact-certified merely by
boundedness or finite-place truncation.

The compactness shortcut is therefore reduced to the following exact analytic
statement:
\[
  P_\infty\tau_{\ell}(I-P_\infty)
  \quad\text{is compact/Hilbert--Schmidt after the legal quotient}
\]
for every finite-place shift \(\ell\) appearing in the local-factor Bohr
spectrum.

If this statement fails for any persistent mode, the noncompact sector is
precisely the sector that must be charged by the Hecke/Dirichlet source
lower-frame ladder.

\section{Membrane status}

In Six Birds membrane language, the finite-place semilocal metric is now a
declared \(P_3/P_4\) route-refinement object, but its compactness is not
accepted.  The honest status is
\[
  \texttt{source\_absorbable\_candidate}
\]
unless the Sonin/prolate off-diagonal compactness theorem is proved.  The
unresolved noncompact block is also an adequacy-risk object: it may contain
directions visible to off-critical-zero dissolving probes but not yet priced by
the native semilocal membrane.  Such directions belong in the \(\Xi\)-ledger
until paid by compact tail or source coercivity.

\end{document}
